% =============================================================================
% Kapitel 7: Fazit und Ausblick
% =============================================================================

\chapter{Fazit und Ausblick}
\label{chap:fazit}

\section{Beantwortung der Forschungsfragen}
\label{sec:ff-beantwortung}

\textbf{FF1} fragt, welche Anforderungen von RFC~9868 sich unter Linux und IPv4 im Benutzerraum
vollständig, teilweise oder nicht erfüllen lassen. Die Zuordnung im geprüften Umfang steht in
\tabref{tab:fazit-synthese}. Die sieben Teilumsetzungen sind Implementierungs- und Entwurfsgrenzen, keine
nachgewiesenen Plattformgrenzen. Der Index selbst mischt markierte BCP-14-Anforderungen, unmarkierte
normative API-Formen, Beschreibungen, Leitlinien und zusammengesetzte Aussagen; er ist ein Arbeitsindex
und kein Verzeichnis aller Pflichten des Standards. Das Ergebnis beantwortet FF1 deshalb auf genau diesem
erklärten Prüfraster und beweist weder die Satzgenauigkeit der Zerlegung noch vollständige Konformität
mit allen Aussagen von RFC~9868.

\textbf{FF2} fragt, in welchem Umfang die Surplus Area auf realen Pfaden bis zur Gegenstelle erhalten
bleibt. Im europäischen Messdreieck und auf den drei transatlantischen Hostpaaren zwischen den
AWS-US-Endpunkten und den europäischen Gegenstellen blieb sie in den jeweiligen Messfenstern erhalten.
Die VMware-NAT wirkte als \textbf{Normalisierer} und entfernte sie. Der gebridgte
Mobilfunk-Hotspot-Pfad und die Google-Fabric wirkten als \textbf{Präsenz-Verwerfer} und verwarfen die
gesamte Klasse in beiden Richtungen; auf einzelnen Hinwegen zu AWS-Regionen in Asien-Pazifik trat
derselbe Totalverlust quellabhängig im Transit auf. Die beobachteten \textbf{Prüfsummen-Gates} verwarfen
Datagramme, deren Einerkomplementsumme über Pseudo-Header und gesamte IP-Nutzlast nicht aufging;
regelkonforme Datagramme mit Null-Pad und gültiger OCS passierten sie. Jede dieser Aussagen gilt nur für
das beobachtete Paar, die Richtung und das Messfenster.

\begin{table}[htbp]
    \centering
    \caption{Synthese der Forschungsfragen im geprüften Umfang}
    \label{tab:fazit-synthese}
    \small
    \begin{tabularx}{\textwidth}{@{}L{1.0cm}L{4.1cm}YL{4.2cm}@{}}
        \toprule
        \textbf{FF} & \textbf{Ergebnis} & \textbf{Beobachtete Grenze} & \textbf{Geltungsbereich} \\
        \midrule
        FF1 & 57 technisch vollständig erfüllbar; 0 teilweise; 0 nicht erfüllbar
        & Referenzimplementierung: 50 vollständig, 7 teilweise; 10 nicht anwendbar oder Leitlinie
        & 67-zeiliger Arbeitsindex; Linux, IPv4, Benutzerraum, Raw Sockets \\
        FF2 & Surplus Area je nach Pfad erhalten, entfernt oder als Klasse verworfen
        & Normalisierer, Prüfsummen-Gates und Präsenz-Verwerfer
        & beobachtetes Paar, Richtung und Messfenster \\
        \bottomrule
    \end{tabularx}
\end{table}

Aus \tabref{tab:fazit-synthese} geht hervor, dass beide Ergebnisse unabhängig voneinander ausfallen. Die
Leitfrage, wie sich UDP um Transportoptionen erweitern lässt, ohne seine grundlegenden Eigenschaften
aufzugeben, ist damit zweigeteilt zu beantworten. Auf der Entwurfsseite bleibt der UDP-Header
unverändert, die Optionen liegen in der Surplus Area, und ein Empfänger liefert die Nutzdaten auch dann
aus, wenn eine SAFE-Option fehlschlägt; Nachrichtenorientierung, Verbindungslosigkeit und der Verzicht
auf Sequenznummern, Fenster und Zeitgeber bleiben erhalten (\secref{subsec:udp-eigenschaften}). Die
einzige belegte Ausnahme ist die Reassemblierung von FRAG-Datagrammen: Sie legt am Empfänger einen
zeitlich und mengenmäßig begrenzten Zustand an, und die Messungen zeigen dessen Verhalten an der
Mengengrenze und am Zeitablauf (\tabref{tab:eval-frag-softstate}). Auf der Pfadseite zeigen die
Messungen, dass die Abwärtskompatibilität des Wire-Formats keine allgemeine Zustellbarkeit garantiert.
Technische Umsetzbarkeit am Endpunkt und Durchlässigkeit des Pfads sind unabhängige Eigenschaften.

\section{Ausblick}
\label{sec:ausblick}

Für eine Weiterentwicklung der Bibliothek liegen die sieben Teilumsetzungen nahe: der fehlende
Geltungsbereich für Optionen vor oder nach der Fragmentierung, eine vom Aufrufer wählbare Mindestlänge
der Sende-API, quellenspezifische Empfangsregeln, die Bereitstellung der empfangenen Parameter auch für
ignorierte oder fehlgeschlagene Optionen und eine klar dokumentierte Trennung des
Reassemblierungsspeichers zwischen Socket-Kontexten. Die hohe Sende-API könnte zusätzlich die Leitlinie
zur paketweisen Fragmentierungssteuerung abbilden. Eine Prüfkette, die semantische Abweichungen ihrer
Auswerter in einen Fehlerstatus überführt, würde den Rückgabewert wieder aussagekräftig machen.

Für künftige Untersuchungen bleiben die Grenzen dieser Arbeit die naheliegenden Ansatzpunkte: eine
unabhängige Implementierung, weitere Zugangsnetze, IPv6, andere Messzeitpunkte und das Zusammenspiel mit
RFC~9869 \cite{rfc9869} könnten Interoperabilität und zeitliche Stabilität der hier berichteten Befunde
einordnen. Wer auf der lokalen 248-Paket-Voranalyse aufbauen will, sollte sie zuvor versiegeln;
portable innere Manifeste würden die äußeren Archivprüfsummen sinnvoll ergänzen.

\section{Schlusswort}
\label{sec:schlusswort}

Die Arbeit belegt, dass die Umsetzung der UDP-Optionsmechanik im Benutzerraum und ihre Zustellbarkeit
auf realen Pfaden getrennte Fragen sind. Eine funktionierende Referenzimplementierung beseitigt die
pfadabhängigen Normalisierer und Verwerfer nicht. Die getrennte Bewertung verhindert, dass ein grüner
Implementierungstest als Netzbeleg oder eine erfolgreiche Pfadmessung als Konformitätsbeleg
missverstanden wird.
